% !TeX program = xelatex
% 本文件汇总本科设计各章节可用的控制/技术路线框图。
% 若插入 main.tex，请在导言区加入：
% \usepackage{tikz}
% \usetikzlibrary{arrows.meta,positioning,calc,fit}

\tikzset{
  ctrlBlock/.style={
    draw,
    rounded corners=2pt,
    align=center,
    minimum height=8mm,
    minimum width=24mm,
    inner xsep=4pt,
    inner ysep=3pt,
    font=\small
  },
  ctrlSmall/.style={
    draw,
    rounded corners=2pt,
    align=center,
    minimum height=7mm,
    minimum width=20mm,
    inner xsep=3pt,
    inner ysep=2pt,
    font=\footnotesize
  },
  ctrlSum/.style={
    draw,
    circle,
    inner sep=0pt,
    minimum size=5.5mm,
    font=\small
  },
  ctrlGroup/.style={
    draw,
    dashed,
    rounded corners=3pt,
    inner sep=5pt
  },
  ctrlArrow/.style={
    -{Latex[length=2.2mm]},
    line width=0.45pt
  },
  ctrlDashArrow/.style={
    -{Latex[length=2.2mm]},
    dashed,
    line width=0.45pt
  }
}

\begin{figure}[htbp]
  \centering
  \resizebox{0.98\textwidth}{!}{%
  \begin{tikzpicture}
    \node[ctrlBlock] (task) at (0,0) {链式自重组\\轮足机器人};
    \node[ctrlBlock] (model) at (3.0,0) {单体/链式\\建模};
    \node[ctrlBlock] (single) at (6.0,0) {单体稳定\\控制};
    \node[ctrlBlock] (coop) at (9.0,0) {链式协同\\控制};
    \node[ctrlBlock] (exp) at (12.0,0) {实验验证\\与对比};

    \node[ctrlSmall] (kin) at (3.0,-1.35) {运动学约束};
    \node[ctrlSmall] (lqr) at (6.0,-1.35) {LQR 基础控制};
    \node[ctrlSmall] (dmpc) at (9.0,-1.35) {DMPC / ADRC};

    \draw[ctrlArrow] (task) -- (model);
    \draw[ctrlArrow] (model) -- (single);
    \draw[ctrlArrow] (single) -- (coop);
    \draw[ctrlArrow] (coop) -- (exp);
    \draw[ctrlArrow] (model) -- (kin);
    \draw[ctrlArrow] (single) -- (lqr);
    \draw[ctrlArrow] (coop) -- (dmpc);
    \draw[ctrlDashArrow] (exp.south) -- ++(0,-2.35) -| node[pos=0.28, below, font=\footnotesize] {性能评价与参数修正} (model.south);
  \end{tikzpicture}
  }%
  \caption{第 1 章研究内容与控制方案总体关系框图}
  \label{fig:ch1_control_route}
\end{figure}

\begin{figure}[htbp]
  \centering
  \resizebox{0.98\textwidth}{!}{%
  \begin{tikzpicture}
    \node[ctrlBlock] (cmd) at (0,0) {头车输入\\$v_1,\omega_1$};
    \node[ctrlBlock] (pose) at (3.0,0) {各单体位姿\\$x_i,y_i,\phi_i$};
    \node[ctrlBlock] (joint) at (6.0,0) {铰接角计算\\$\theta_i$};
    \node[ctrlBlock] (constraint) at (9.0,0) {铰接点速度\\一致约束};
    \node[ctrlBlock] (prop) at (12.0,0) {后车期望输入\\$v_{i+1},\omega_{i+1}$};

    \node[ctrlSmall] (leg) at (3.0,-1.45) {腿部运动学\\$q_j \rightarrow L_0,\Phi_0$};
    \node[ctrlSmall] (seq) at (9.0,-1.45) {逐级传播\\形成链式期望序列};

    \draw[ctrlArrow] (cmd) -- (pose);
    \draw[ctrlArrow] (pose) -- (joint);
    \draw[ctrlArrow] (joint) -- (constraint);
    \draw[ctrlArrow] (constraint) -- (prop);
    \draw[ctrlArrow] (pose) -- (leg);
    \draw[ctrlArrow] (prop) |- (seq);
    \draw[ctrlDashArrow] (seq.south) -- ++(0,-0.7) -| node[pos=0.24, below, font=\footnotesize] {更新下一单体} ($(leg.west)+(-0.35,-0.05)$) |- (pose.west);
  \end{tikzpicture}
  }%
  \caption{第 2 章链式系统运动学建模框图}
  \label{fig:ch2_chain_kinematics}
\end{figure}

\begin{figure}[htbp]
  \centering
  \resizebox{0.98\textwidth}{!}{%
  \begin{tikzpicture}
    \node[ctrlBlock] (ref) at (0,0) {参考量\\$r=[v_x^{ref},\omega_z^{ref}]^T$};
    \node[ctrlBlock] (map) at (3.5,0) {参考状态映射\\$R_{ref}r$};
    \node[ctrlSum] (sum) at (5.8,0) {$+$};
    \node[ctrlBlock] (lqr) at (8.2,0) {LQR 状态反馈\\$u=-Kx_{err}$};
    \node[ctrlBlock] (torque) at (11.7,0) {力矩分配\\轮端/腿部/尾部};
    \node[ctrlBlock] (plant) at (15.2,0) {线性化单体模型\\$\dot{x}=Ax+Bu$};
    \node[ctrlBlock] (out) at (18.7,0) {输出\\$y=[v_x,\omega_z]^T$};

    \node[ctrlSmall] (fit) at (8.2,-1.45) {腿长参数化增益\\$K(l_l,l_r)$};

    \draw[ctrlArrow] (ref) -- (map);
    \draw[ctrlArrow] (map) -- node[above, font=\footnotesize] {$+$} (sum);
    \draw[ctrlArrow] (sum) -- (lqr);
    \draw[ctrlArrow] (torque) -- (plant);
    \draw[ctrlArrow] (plant) -- (out);
    \draw[ctrlArrow] (fit) -- (lqr);
    \draw[ctrlArrow] (lqr) -- (torque);
    \draw[ctrlArrow] (out.south) -- ++(0,-2.45) -| node[pos=0.22, below, font=\footnotesize] {状态反馈 $-x$} (sum.south);
  \end{tikzpicture}
  }%
  \caption{第 3 章单体 LQR 基础控制闭环框图}
  \label{fig:ch3_lqr_closed_loop}
\end{figure}

\begin{figure}[htbp]
  \centering
  \resizebox{0.98\textwidth}{!}{%
  \begin{tikzpicture}
    \node[ctrlBlock] (traj) at (0,0) {链式期望轨迹\\$r_i$};
    \node[ctrlBlock] (measure) at (3.5,0) {本车及邻车状态\\$z_{i-1},z_i,z_{i+1}$};
    \node[ctrlBlock] (dmpc) at (7.0,0) {局部 DMPC\\滚动优化};
    \node[ctrlBlock] (cmd) at (10.4,0) {协同层输出\\$u_i^\ast=[v_i,\omega_i]^T$};
    \node[ctrlBlock] (inner) at (14.1,0) {下层单体控制器\\LQR / LQR-ADRC};
    \node[ctrlBlock] (robot) at (18.0,0) {第 $i$ 个机器人\\运动响应};

    \node[ctrlSmall] (comm) at (3.5,-1.5) {邻车预测序列\\信息交换};
    \node[ctrlSmall] (cost) at (7.0,-1.5) {跟踪误差\\铰接约束\\输入变化率};

    \draw[ctrlArrow] (traj) -- (measure);
    \draw[ctrlArrow] (measure) -- (dmpc);
    \draw[ctrlArrow] (dmpc) -- (cmd);
    \draw[ctrlArrow] (cmd) -- (inner);
    \draw[ctrlArrow] (inner) -- (robot);
    \draw[ctrlArrow] (cost) -- (dmpc);
    \draw[ctrlArrow] (comm) -- (measure);
    \draw[ctrlDashArrow] (robot.south) -- ++(0,-2.45) -| node[pos=0.23, below, font=\footnotesize] {状态更新} (measure.south);
  \end{tikzpicture}
  }%
  \caption{第 4 章分层 DMPC 协同控制框图}
  \label{fig:ch4_dmpc_hierarchical}
\end{figure}

\begin{figure}[htbp]
  \centering
  \resizebox{0.98\textwidth}{!}{%
  \begin{tikzpicture}
    \node[ctrlBlock] (scheme) at (0,0) {控制方案\\LQR / DMPC / LQR-ADRC};
    \node[ctrlBlock] (platform) at (4.2,0) {实验平台\\单体/链式系统};
    \node[ctrlBlock] (sensors) at (8.3,0) {数据采集\\姿态/速度/腿长/铰接角};
    \node[ctrlBlock] (metrics) at (12.6,0) {评价指标\\误差/峰值/稳定时间};
    \node[ctrlBlock] (compare) at (16.7,0) {策略对比\\性能分析};

    \node[ctrlSmall] (dist) at (4.2,-1.55) {速度指令\\地形/外力扰动};
    \node[ctrlSmall] (log) at (12.6,-1.55) {离线统计\\曲线分析};

    \draw[ctrlArrow] (scheme) -- (platform);
    \draw[ctrlArrow] (platform) -- (sensors);
    \draw[ctrlArrow] (sensors) -- (metrics);
    \draw[ctrlArrow] (metrics) -- (compare);
    \draw[ctrlArrow] (dist) -- (platform);
    \draw[ctrlArrow] (sensors) -- (log);
    \draw[ctrlArrow] (log) -- (metrics);
    \draw[ctrlDashArrow] (compare.south) -- ++(0,-2.55) -| node[pos=0.25, below, font=\footnotesize] {参数整定与方案修正} (scheme.south);
  \end{tikzpicture}
  }%
  \caption{第 5 章实验测试与结果分析流程框图}
  \label{fig:ch5_experiment_evaluation}
\end{figure}
